\section{Kontextfreie Sprachen -- Teil 2}


\subsection{Deterministisch kontextfreie Sprachen}

\begin{Def}
Wir definieren einen \emph{Kellerautomaten mit akzeptierenden Zuständen} als 7-Tupel
\[
	\mathcal{E} = (\Sigma,Q,\Gamma,q^\text{init},\Zinit,\delta,F)
\]
mit den Komponenten $\Sigma,Q,\Gamma,q^\text{init},\Zinit,\delta$ wie für einen gewöhnlichen Kellerautomaten und $F \subseteq Q$ als die Menge der akzeptierenden Zustände.

\smallskip

Weiterhin definieren wir das Akzeptanzverhalten folgendermaßen: \\
Eine Konfiguration $(q, w, \gamma)$ eines Kellerautomaten mit akzeptierenden Zuständen ist \emph{akzeptierend}, falls $q \in F$ und $w = \varepsilon$.
Die von einem Kellerautomaten mit akzeptierenden Zuständen $\mathcal{E}$ \emph{akzeptierte Sprache} ist
\[
	L(\mathcal{E}) = \{w \in \Sigma^* \mid \exists q \in F, \gamma \in \Gamma^*: (q^\text{init}, w, Z^\text{init}) \ConfRel^* (q, \varepsilon, \gamma)\}.
\]
Hierbei ist $\ConfRel$ die Schrittrelation genau wie für gewöhnliche Kellerautomaten.
\end{Def}

\begin{lemma}
Zu jedem (gewöhnlichen) Kellerautomaten gibt es einen Kellerautomaten mit akzeptierenden Zuständen, der die gleiche Sprache akzeptiert.
Gleiches gilt in die andere Richtung.
\end{lemma}
% \begin{proof}
% Siehe Übungsblatt~8 Aufgabe~5.
% \end{proof}




\begin{Def}[name={[DPDA]}]
Wir nennen sowohl einen gewöhnlichen Kellerautomaten als auch einen Kellerautomaten mit akzeptierenden Zuständen \emph{deterministisch},
wenn für alle $q\in Q,a\in\Sigma,Z\in\Gamma$ gilt, dass $|\delta(q,a,Z)| + |\delta(q,\Eps,Z)| \leq 1$.

Wir verwenden die Abkürzung \ac{DPDA} für den deterministischen Kellerautomaten mit akzeptierenden Zuständen.
\end{Def}

\begin{lemma}[name={[\acs*{DPDA}, der gesamte Eingabe verarbeitet]}]
        \label{lem:DPDA ges. Eingabe}
        Zu jedem \ac{DPDA} $\mathcal{D}$ gibt es einen äquivalenten \ac{DPDA} $\mathcal{D}_e$, der
        jede Eingabe bis zum Ende liest.
\end{lemma}
% \draftnote{23.12.16}
\begin{proof}

  Sei $\mathcal{D} = (\Sigma, Q, \Gamma, \qinit, \Zinit, \delta, F)$. 
  Es gibt zwei Möglichkeiten, warum $\mathcal{D}$ nicht die gesamte Eingabe verarbeitet:
  Die Transitionsrelation ist nicht total 
%   \footnote{
%   d.h. es gibt Konfigurationen mit nicht-leerem Keller
%   so dass $a \not \Rightarrow^* b$} 
  oder der Automat bleibt bei
  leerem Keller stecken.

  Abhilfe: Führe einen Senkzustand ein, auf dem $\mathcal{D}_S$ weiterrechnet, wenn $\mathcal{D}$ nicht mehr weiterrechnen kann. 
  Definiere dazu $\mathcal{D}_S = (\Sigma, Q_S, \Gamma_S, \qinit_S, \Zinit_S, \delta_S, F_S)$ mit
  \begin{align*}
  Q_S &= Q \cup \{\qinit_S, q_s\} \text{ mit neuen Zuständen: $\qinit_S$ neuer Startzustand und $q_s$ Senkzustand} \\
  F_S &= F \\
  \Gamma_S &= \Gamma \cup \{\Zinit_S\} \text{ neues Kellerbodensymbol}
  \end{align*}
  und Transitionsfunktion $\delta_S$ gegeben durch
  \[\delta_S (\qinit_S, \varepsilon, \Zinit_S) = \{ (\qinit, \Zinit\Zinit_S) \} \]
  \begin{itemize}
  \item Transitionsfunktion totalisieren: Für alle $q\in Q$, $Z \in \Gamma$
    \begin{align*}
      \delta_S (q, \varepsilon, Z) &=
      \begin{cases}
        \delta (q, \varepsilon, Z) & \text{falls
        }\delta (q, \Eps, Z)\ne\emptyset\vee\exists a\in\Sigma, \delta (q, a, Z)\ne \emptyset \\
        \{ (q_s, Z) \} & \text{sonst}
      \end{cases} \\
      \delta_S (q, a, Z) &=
      \begin{cases}
        \delta (q, a, Z) & \text{falls }\delta (q, \Eps, Z)\ne\emptyset\vee\delta (q,
        a, Z) \ne \emptyset \\
        \{ (q_s, Z) \} & \text{sonst}
      \end{cases}
    \end{align*}
    Intuition: Wenn $\delta(q, a, Z) = \emptyset$ und es keinen $\epsilon$-Übergang gibt, kann das Wort von $\mathcal{D}$ nicht weiter abgearbeitet werden. $\mathcal{D}_S$ geht nun in Senkzustand $q_s$ und arbeitet dort weiter. Analog für $\delta(q, \epsilon, Z)$.
  \item Kellerunterlauf vermeiden: $\mathcal{D}$ hat $\Zinit$ abgeräumt, so dass $\Zinit_S$
    sichtbar geworden ist. Füge eine $\epsilon$-Transition in
     Senkzustand $q_s$ hinzu, indem man für alle $q\in Q$ definiert:
    \[\delta_S (q, \varepsilon, \Zinit_S) = \{ (q_s, \Zinit_S) \} \]
  In $q_s$ kann $\mathcal{D}_S$ nun das Restwort für alle $a\in\Sigma$, $Z\in \Gamma_S$ abarbeiten:
    \[\delta_S (q_s, a, Z) = \{ (q_s, Z) \}\]
    Bemerkung: $q_s \not \in F_S$, da ein Kellerunterlauf
    das Akzeptieren des eingelesen Worts ausschließt.
  \end{itemize}

  Es verbleibt zu zeigen:
  \begin{itemize}
   \item $\forall w\in \Sigma^*$ $\exists q\in Q_S$,
  $\gamma\in\Gamma_S^*: (\qinit_S, w, \Zinit_S) \ConfRel^* (q,
  \varepsilon, \gamma)$.
  
  (Induktion über die Länge von $w$)
  \item $L (\mathcal{D}) = L (\mathcal{D}_S)$.
  \qedhere
  \end{itemize}
\end{proof}

\begin{Def}
 Wir nennen eine Sprache $L$ \emph{deterministisch kontextfrei} wenn es einen \ac{DPDA} gibt der $L$ akzeptiert.
\end{Def}


\begin{Satz}[name={[Abgeschlossenheit der deterministischen \acs*{CFL}]}]
        Die deterministisch kontextfreien Sprachen sind unter Komplement abgeschlossen.
\end{Satz}
\begin{proof}
  Sei $L$ deterministisch kontextfrei. 
  Nach \autoref{lem:DPDA ges. Eingabe} gibt es einen \ac{DPDA} $\mathcal{D}$ mit $L(\mathcal{D})=L$, der die komplette Eingabe liest.

  Konstruiere einen \ac{DPDA} $\mathcal{D}_K$, sodass $ L(\mathcal{D}_K) = \overline{L}$. 
  Definiere dazu $Q_K= Q \times \{0,1,2\}$. Bedeutung des zusätzlichen "`Flags"' in $\{0,1,2\}$:
  \begin{itemize}
  \item $0$: seit dem Lesen des letzten Symbols $a \in \Sigma$ wurde kein akzeptierender Zustand durchlaufen
  \item $1$: seit dem Lesen des letzten Symbols $a \in \Sigma$ wurde mindestens ein
    akzeptierender Zustand durchlaufen 
  \item $2$: markiert einen akzeptierenden Zustand in $\mathcal{D}_K$
  \end{itemize}
        
  $F_K=Q\x \{2\}$

  Definiere Hilfsfunktion $\mathit{acc}:Q\to \{0,1\}$ durch 
  \begin{displaymath}
    \mathit{acc} (q) =
    \begin{cases}
      0 & q\notin F \\ 1 & q \in F
    \end{cases}
  \end{displaymath}
  \begin{align*}
    \qinit_K &= (\qinit,\mathit{acc} (\qinit))
  \end{align*}
  Wir definieren $\delta'$ für alle $q \in Q$, $Z\in\Gamma$ folgendermaßen.
  
  Falls $\delta(q,\Eps,Z) = \{(q',\gamma)\}$, dann
  \begin{align*}
    \delta'((q,0),\Eps,Z) &= \{(q',\mathit{acc} (q'),\gamma)\}
    \\
    \delta'((q,1),\Eps,Z) &= \{(q',1,\gamma)\}
  \end{align*}
  Falls $\delta(q,a,Z) = \{(q',\gamma)\}$, dann
  \begin{align*}
    \delta'((q,0), \Eps, Z) &= \{ ((q,2), Z) \} \\
    \delta'((q,2), a, Z) &=\{(q',\mathit{acc} (q'), \gamma)\}
    \\
    \delta'((q,1), a, Z) &=
    \{(q',\mathit{acc} (q'), \gamma)\}
    \qedhere    %TODO Intuition?
  \end{align*}
\end{proof}

\begin{Satz}
Die deterministisch kontextfreien Sprachen sind eine echte Teilmenge der kontextfreien Sprachen.
\end{Satz}

\begin{proof}
 Offensichtlich gibt es eine Teilmengenbeziehung, denn die von nichtdeterministischen Kellerautomaten mit akzeptierendem Zustand akzeptierten Sprachen sind gerade die kontextfreien Sprachen und jeder deterministische Kellerautomat ist ein spezieller nichtdeterminischer Kellerautomaten mit akzeptierenden Zuständen.
 Die deterministisch kontextfreien Sprachen sind unter Komplement abgeschlossen, die kontextfreien Sprachen sind es nicht.
 Die beiden Mengen von Sprachen können also nicht gleich sein.
\end{proof}



\hide{

\begin{Satz}
    Die deterministischen \ac{CFL} sind \textbf{nicht} unter Vereinigung und Durchschnitt abgeschlossen.
\end{Satz}
\begin{proof}
    Betrachte
    \begin{align*}
        L_1 &= \{ a^nb^nc^m \mid n, m \ge 1 \} \\
        L_2 &= \{ a^mb^nc^n \mid n, m \ge 1 \}
    \end{align*}
    Sowohl $L_1$ als auch $L_2$ sind DCFL, aber $L_1 \cap L_2 = \{ a^nb^nc^n \mid n \ge 1\}$ ist nicht kontextfrei.
    
    DCFL ist nicht abgeschlossen unter Vereinigung. Angenommen doch: Seien $U, V$ DCFL. Dann sind auch $\overline{U}$ und $\overline{V}$ DCFL. Bei Abschluss unter Vereinigung wäre $\overline{U} \cup \overline{V}$ eine DCFL und somit auch $\overline{\overline{U} \cup \overline{V}} = U \cap V$, ein Widerspruch gegen den ersten Teil.
\end{proof}
\begin{Satz}
    DCFL ist abgeschlossen unter Schnitt mit REG.
\end{Satz}
\begin{proof}
    Sei $L$ DCFL und $R$ regulär.
    Konstruiere das Produkt aus einem DPDA für $L$ und einem DFA für $R$.
    Offenbar ist das Ergebnis ein DPDA, der $L\cap R$ akzeptiert. \footnote{
    Intution des Beweis: Konstruiere Produktautomaten, der akzeptiert gdw. der DFA und NDPA akzeptieren. Idee: Definiere Tupel $(q_1, q_2)$, DPDA arbeitet auf $q_1$, DFA arbeitet auf $q_2$. Beispielsweise werden $\epsilon$-Übergänge des DPDA nur auf der "`DPDA-Seite"' d.h. auf $q_1$ abgearbeitet, ohne Beeinflussung von $q_2$.
    }

    $L = L (M_1)$ mit $M_1 = (Q_1, \Sigma, \Gamma, q_{01}, \Zinit,
    \delta_1, F_1)$ \ac{DPDA}

    $R = L (M_2)$ mit $M_2 = (Q_2, \Sigma, q_{02}, \delta_2, F_2)$
    \ac{DFA}

    Konstruiere $M'= (Q, \Sigma, \Gamma, \qinit, \Zinit, \delta, F)$ mit
    \begin{itemize}
    \item $Q = Q_1 \times Q_2$
    \item $\qinit = (q_{01}, q_{02})$
    \item $F = F_1  \times F_2$
    \item Falls $\delta_1 (q_1, \varepsilon, Z) = \{(q_1', \gamma)\}$, dann gilt für alle $q_2\in Q_2$: 
      \[ \delta ((q_1, q_2), \varepsilon, Z)
      = \{((q_1', q_2), \gamma)\} \text{.     (Ansonsten leer)} \]
    \item Falls $\delta_1 (q_1, a, Z) = \{(q_1', \gamma)\}$, dann gilt für alle $q_2\in Q_2$: 
      \[ \delta ((q_1, q_2), a, Z) =
      \{((q_1', \delta_2 (q_2, a)), \gamma)\} \text{   (Ansonsten leer)} \]
    \end{itemize}
    Zu zeigen ist noch $L (M') = L (M_1) \cap L (M_2)$.
\end{proof}
\begin{Satz}
    Sei $L$ DCFL und $R$ regulär.
    Es ist entscheidbar, ob $R=L$, $R\subseteq L$ und $L=\Sigma^*$.
\end{Satz}
\begin{proof}
    Es gilt $R\subseteq L$ gdw.\ $R \cap \overline{L} = \emptyset$.
    
    Weiter ist $R = L$ gdw.\ $R\subseteq L$ und $L \subseteq R$. Für den zweiten Teil betrachte $L\cap \overline{R}$.
    
    Für kontextfreie Sprachen ist $L\ne \emptyset$ entscheidbar, also betrachte $L=\Sigma^*$ gdw.\ $\overline{L}=\emptyset$.
\end{proof}















\begin{Satz}\label{satz:5.1}
  \begin{align*}
                L \in \ac{CFL}  \text{ gdw } L =L(M) \text{ für einen \ac{NPDA} $M$}
  \end{align*}
\end{Satz}
\begin{proof}\hfill
        \begin{itemize}
        \item \ac{CFG} zu \ac{NPDA}:


  \item \ac{NPDA} zu \ac{CFG}:

    Zunächst zeigen wir, dass es genügt, \ac{NPDA}s zu betrachten, die bei jeder Transition Wörter der maximalen Länge $2$ auf den Keller schreiben:



\end{itemize}
  
\end{proof}



\begin{Satz} \textbf{DPDA Äquivalenzproblem}
    Seien $L_1, L_2$ DCFL. Dann ist $L_1 = L_2$ entscheidbar.
\end{Satz}

\begin{proof}
    Siehe Senizergues (2000) und Stirling (2001).
\end{proof}

Wir betrachten zum Ende des Kapitels noch eine praktische Fragestellung: Wie sieht man einer CFL an, dass sie von einem DPDA akzeptierbar ist?

Sei $\mathcal{G} = (\Sigma, N, P, S)$.
Wir können nach Satz \ref{satz:5.1} einen \ac{PDA} für $\mathcal{G}$ konstruieren, mit
\begin{align*}
  \delta(q, a, a) &= \{(q, \Eps)\} \quad a \in \Sigma \\
  \delta(q, \Eps, A) &= \{(q, \beta)  \mid A \to \beta \in P\}
\end{align*}
Die Transitionen für Eingabezeichen $a \in \Sigma$ sind deterministisch.
Die $\Eps$-Transitionen sind es nicht unbedingt.

Die Idee ist nun, den Automaten mit einem Symbol Lookahead \footnote{Das Problem ist, dass ein deterministischer \ac{PDA} nicht "`alle mögl. Regeln gleichzeitig ausprobieren"' kann, so wie ein NDPA. Möchte der Automat ein Eingabesymbol $a$ matchen und kann aktuell mehrere Produktionen anwenden, müsste er in die Zukunft sehen können, um diejenige zu wählen, die das gewünschte $a$ an erster Position erzeugt. Diesen "`Blick in die Zukunft"' gewähren die $first(A)$ Mengen: Sie geben für jedes Nichtterminalsymbol $A$ an, welche Terminale $a \in \Sigma$ als Präfix von $yield(A)$ auftreten können.} 
erweitern. Dieses Symbol wird jeweils im Zustand des Automaten
gespeichert. Der \ac{PDA} für $L (\mathcal{G})$ mit Lookahead hat nun
folgende Komponenten:
\begin{itemize}
\item $Q = \Sigma\cup \{\varepsilon, \$\}$
\item $\qinit = [\Eps]$ (zu Beginn ist der Lookahead leer)
\item $F = [\$]$  (das Symbol $\$$ markiert das Ende der Eingabe)
\end{itemize}
und die Transitionsfunktion
\begin{align*}
  \delta ([\Eps], a, Z) & = \{([a], Z) \} & & \text{lade Lookahead} \\
  \delta ([a], \Eps, a) & = \{([\Eps], \varepsilon) \} && \text{match}
  \\
  \delta ([a], \Eps, A) & = \{([a], \beta) \mid A\to\beta\in P, a \in
  \mathit{first} (\beta) \} && \text{select}
\end{align*}
Der Automat startet in der Konfiguration $([\Eps], w\$, S\$)$.

\textbf{Beispiel}: Sei ein \ac{CFG} gegeben mit den Produktionen $S \to ( S )$ $ |$ $ a$. Für die Eingabe $(a)$ ergibt sich folgende Abarbeitung:
\begin{align*}
              & ([\Eps], (a), S\$ )   && \text{Lade Lookahead ,('}\\
  \Rightarrow & ([(], (a), S\$ )   && \text{Select: Da ,(' $\in first($ $(a)$ $ )$, wende $S \to (a)$ an}\\
  \Rightarrow & ([(], (a), (S)\$ ) && \text{Match ,('}\\
  \Rightarrow & ([\Eps], a), S)\$) \\
  \Rightarrow & ...
\end{align*}

\begin{Def}
  Sei $\mathcal{G} = (\Sigma, N, P, S)$ \ac{CFG} und $\beta \in
  (N\cup\Sigma)^*$
  \begin{align*}
    \mathit{first} (\beta) &= \{ a \in \Sigma \mid \exists w\in
    \Sigma^* , \beta \=>^* aw \} \cup \{\Eps \mid \beta \=>^* \Eps \}
  \end{align*}
\end{Def}

Spezifikation von $\mathit{first}(\beta)$
\begin{align*}
  \mathit{first} (\Eps) & = \{ \Eps \} \\
  \mathit{first} (a\beta) & = \{ a \} \\
  \mathit{first} (A\beta) & =
  \begin{cases}
    \mathit{first} (A) & \Eps \notin \mathit{first (A)} \\
    \mathit{first} (A) \setminus \{\Eps\} \cup \mathit{first} (\beta)
    & A\=>^* \Eps
  \end{cases} \\
  \mathit{first} (A) &= \bigcup  \{ \mathit{first} (\beta) \mid A \to
  \beta \in P \}
\end{align*}

\textbf{Algorithmus} zur Berechnung von $\mathit{first} (A)$ für alle $A\in N$

Sei $FI[A] \subseteq N$ ein Feld indiziert mit Nichtterminalsymbolen.

\begin{itemize}
\item[] For each $A \in N$: $FI[A] \gets \emptyset$
\item[] Repeat
  \begin{itemize}
  \item[] For each $A\in N$: $FI'[A] \gets FI[A]$
  \item[] For each $A \to \beta \in P$
    \begin{itemize}
    \item[] $FI[A] \gets FI[A] \cup \mathit{first}_{FI} (\beta)$
    \end{itemize}
  \end{itemize}
\item[] Until $\forall A\in N$: $FI'[A] = FI[A]$
\end{itemize}

Dabei ist
\begin{align*}
  \mathit{first}_{FI} (\Eps) & = \{ \Eps \} \\
  \mathit{first}_{FI} (a\beta) & = \{ a \} \\
  \mathit{first}_{FI} (A\beta) & =
  \begin{cases}
    FI[A] & \Eps \notin FI[A] \\
    FI[A] \setminus \{\Eps\} \cup \mathit{first}_{FI} (\beta)
    & \Eps\in FI[A]
  \end{cases} 
\end{align*}

\textbf{Beispiel}: Betrachte eine  Grammatik arithmetische Ausdrücke
mit Startsymbol $S$ und $N = \{E, T, F\}$, $\Sigma = \{a, -, *\}$ und Produktionen
\begin{align*}
  E & \to T E' & E' & \to -T E' \mid \varepsilon \\
  T & \to F T' & T' & \to *F T' \mid \varepsilon \\
  F & \to a
\end{align*}

Tabelle der Werte von $FI[A]$ wobei Zeile $i$ den Werten in $FI$ nach
dem $i$-ten Schleifendurchlauf entspricht.

\begin{displaymath}
  \begin{array}{c||c|c|c|c|c|}
    FI & E & E' & T & T' & F \\\hline
    0  & \emptyset & \emptyset & \emptyset & \emptyset & \emptyset \\
    1  & \emptyset & \{ -, \Eps \} & \emptyset & \{ *, \Eps \} & , \{a\} \\
    2  & \emptyset & \{ -, \Eps \} & \{a\} & \{ *, \Eps \} & , \{a\} \\
    3  & \{a\} & \{ -, \Eps \} & \{a\} & \{ *, \Eps \} & , \{a\} \\
    4  & \{a\} & \{ -, \Eps \} & \{a\} & \{ *, \Eps \} & , \{a\} \\
  \end{array}
\end{displaymath}

Ergebnis nach vier Durchläufen. 

Anmerkung: first ist nicht die vollständige Lösung des Problems. Für
den Lookahead müssen auch noch die Symbole betrachtet werden, die
\emph{nach} einem bestimmten Nichtterminalsymbol auftreten können. Mehr dazu
in Vorlesung Compilerbau. 
}

%%% Local Variables:
%%% mode: latex
%%% TeX-master: "Info_3_Skript_SS2026"
%%% End:(\dots)
